\documentclass[12pt, letterpaper]{article}
\usepackage[margin=1.2in]{geometry}
\usepackage{ebgaramond}
\usepackage[T1]{fontenc}
\usepackage{microtype}
\usepackage{setspace}
\usepackage{parskip}
\usepackage{titling}
\usepackage{fancyhdr}

\setlength{\parindent}{2em}
\onehalfspacing

\pagestyle{fancy}
\fancyhf{}
\fancyhead[L]{\small\textit{Four Inches}}
\fancyhead[R]{\small\textit{NASA Mission Series v1.16}}
\fancyfoot[C]{\small\thepage}
\renewcommand{\headrulewidth}{0.4pt}

\title{\Huge\textbf{Four Inches}\\[0.5em]\large\textit{A short story from Mission 1.16 --- Mercury-Redstone 1}}
\author{NASA Mission Series\\[0.3em]\small Miles Tiberius Foxglove --- Mukilteo, WA}
\date{March 30, 2026}

\begin{document}
\maketitle
\thispagestyle{empty}

\vspace{1em}
\noindent\rule{\textwidth}{0.4pt}
\vspace{1em}

\noindent The blockhouse at Launch Complex~5 was three hundred feet from the pad, which was close enough to feel the rocket in your teeth when it lit but far enough that you could pretend the concrete walls would save you if it didn't go well. I had my hand on the edge of the console, not because I needed steadying but because touching something solid made the abstraction of telemetry feel less abstract. Numbers on a screen. Voltages. Pressures. Sequencer states. The rocket was out there, tall and white and red in the November Florida sun, and in here it was a column of numbers.

``T minus ten seconds,'' the count came, and I stopped breathing, which is what you do, which is what everybody does, even the ones who've done it twenty times.

I had done it four times. Four countdowns. Two had ended in scrubs. One had gone beautifully --- a Redstone test flight that arced over the Atlantic like God had thrown it. One had ended in a hold at T-minus-forty that stretched into a scrub when a valve actuator showed intermittent resistance. That had been my valve. My telemetry channel. I'd spent six hours with the contractor tracing the wiring before we found a cold solder joint the size of a pinhead.

This was number five. Mercury-Redstone~1. The first Mercury capsule on a Redstone booster. No one aboard --- this was an unmanned test --- but the capsule was real, the escape tower was real, the abort sensing system was real. Everything was real except the astronaut.

\begin{center}
$\bullet\quad\bullet\quad\bullet$
\end{center}

The A-7 lit at T-zero and the blockhouse shook. Not a metaphor. The consoles vibrated. The coffee in Gene's mug rippled. You could feel it in your sternum, a low-frequency pressure wave that bypassed your ears and went straight to your skeleton. The engine was producing seventy-eight thousand pounds of thrust three hundred feet away. The sound arrived a fraction of a second after the light, because sound is slow and light is fast, and there was a tiny interval --- maybe a quarter second --- where you could see the flame but not hear it, and your brain, which evolved in forests where silence after a flash meant lightning, told you something was wrong even though nothing was wrong.

Then the sound caught up and everything was noise.

``Thrust is building,'' someone said. The numbers confirmed it. The A-7 was ramping up through fifty thousand, sixty thousand, seventy thousand pounds of thrust. The holddown clamps were still gripping. The rocket was straining against four steel fingers while sixty-six thousand pounds of vehicle shuddered on the pad.

``Clamps release.''

I watched the altitude readout. It started to move. Slowly --- at first just fractions of a foot per second. The rocket was rising, a column of white and red separating from the pad by inches, the gap between fin and concrete growing like a door opening onto something vast.

Four inches.

\begin{center}
$\bullet\quad\bullet\quad\bullet$
\end{center}

The engine stopped.

There was no warning. No anomalous telemetry, no parameter drifting out of range, no bearing frequency climbing into the danger zone the way the Pioneer~0 turbopump had announced itself two years earlier. The A-7 was running and then it was not. The thrust readout dropped from seventy-eight thousand to zero in the time it takes to blink. The sound stopped. Not faded --- stopped. The silence was louder than the engine.

Nobody spoke. I watched the altitude readout. It showed four inches. Then three. Then two. Then one. Then zero. The rocket had risen four inches off the pad and settled back down on its fins like a man who stands up from a chair, changes his mind, and sits back down.

``What happened?'' someone asked, but nobody answered because nobody knew.

\begin{center}
$\bullet\quad\bullet\quad\bullet$
\end{center}

Then the escape tower fired.

The solid rocket motor in the Launch Escape System ignited with a crack that punched through the blockhouse walls. Through the periscope I watched the tower --- a lattice of steel struts topped by a solid motor --- lift off the capsule and arc away trailing orange flame and white smoke. It flew beautifully. It was, in that moment, the most successful component on the entire vehicle. It had detected loss of thrust, interpreted it as a booster failure, and fired to pull the capsule clear.

Except the capsule was still sitting on top of the rocket. The clamp ring hadn't received a separation command. The tower flew away alone, a rescue system rescuing nothing, and landed on the beach four hundred yards downrange.

Then the drogue parachute deployed.

The capsule's sequencer, now running through its reentry timeline, checked the barometric altitude. Sea level. Below ten thousand feet. Deploy drogue. The small stabilization chute popped out of the top of the capsule and caught the Cape Canaveral breeze. It fluttered over the nose of the capsule like a flag on a very short pole.

Then the main parachute deployed.

The big one. Sixty-three feet of orange-and-white ringsail canopy, designed to lower a Mercury capsule and its astronaut gently into the Atlantic Ocean from an altitude of ten thousand feet. It billowed out in the onshore wind and settled over the rocket like a bedsheet over a piece of furniture in a house someone is leaving.

I stared through the periscope at a fully-fueled Mercury-Redstone rocket standing on Launch Complex~5 draped in its own parachute, and I understood for the first time what Voltaire meant when he said this is the best of all possible worlds.

\begin{center}
$\bullet\quad\bullet\quad\bullet$
\end{center}

We couldn't approach. The rocket had live pyrotechnics --- destruct charges, retrorocket igniters, parachute mortars that may not have fully fired. The propellant tanks were full. Twenty-four thousand kilograms of liquid oxygen and ethanol sitting on the pad with enough chemical energy to flatten the blockhouse. The silver-zinc batteries were still powering every circuit in the vehicle, and any of those circuits could fire any of those pyrotechnics at any moment for any reason or no reason at all.

Chris Kraft made the call. Wait. Do nothing. Let the batteries die. Come back tomorrow.

Twenty-four hours. We sat in the blockhouse for the first six and then moved to the ready room because the blockhouse wasn't designed for sleeping and the coffee had run out. The rocket stood on the pad all night, the parachute billowing gently in the Atlantic breeze, the red aircraft warning light on the capsule blinking in the darkness, a lighthouse warning ships away from a shore that had a rocket on it wearing a parachute.

\begin{center}
$\bullet\quad\bullet\quad\bullet$
\end{center}

I thought about \textit{Dictyostelium discoideum} that night, which is not what you'd expect an engineer to think about at three in the morning while waiting for a rocket's batteries to die. But I'd been reading about them --- the social amoebae, the cellular slime molds. Single-celled organisms that live independently in soil, eating bacteria, minding their own business, until food runs out. Then they do something remarkable. They aggregate. Thousands of individual cells stream toward a central point, following trails of cyclic AMP, and form a multicellular slug. The slug crawls. It behaves as a single organism even though it's thousands of organisms. It rises up into a stalk topped by a fruiting body, and the cells at the top become spores and disperse, and the cells that formed the stalk die.

The slug rises. Then it collapses. The stalk cells sacrifice themselves so the spore cells can travel. The organism that rises is not the organism that disperses. Something goes up. Something else comes down. The rising and the falling are both part of the design.

MR-1 rose four inches. Then it came down. The escape tower rose four hundred feet. Then it came down. The parachute rose sixty-three feet of canopy into the air. Then it came down. Everything went up and everything came down and every single system worked perfectly. The engine started correctly. The abort system detected the failure correctly. The sequencer deployed the parachutes correctly. The 24-hour wait was the correct safety decision. Everything worked. Nothing succeeded.

\begin{center}
$\bullet\quad\bullet\quad\bullet$
\end{center}

By morning, the telemetry team had traced the fault. The tail plug connector at the base of the rocket had two sets of pins --- power and control. When the rocket rose, the plug pulled free. The design assumed a rise of several feet before separation, with the power pins disconnecting first (longer pins, they hold on longer) and the control pins disconnecting second. At four inches, the connector separated but the pins came out in the wrong order. The control pins --- slightly shorter --- disconnected first. For a fraction of a second, ground power was still flowing through the connector while the control-inhibit circuit was open.

A sneak circuit. An unintended path through the partially-connected plug. Ground power flowed through the open inhibit path and reached the engine shutdown sequencer. The sequencer received what looked like a legitimate shutdown command and obeyed it.

The engine did exactly what it was told. The problem was who was doing the telling.

\begin{center}
$\bullet\quad\bullet\quad\bullet$
\end{center}

The fix was elegant in the way that only fixes for stupid problems can be elegant: make the power pins physically longer. By a quarter inch. So that regardless of rise height --- four inches, four feet, four millimeters --- the power pins always disconnect first. The correct sequence is guaranteed by geometry, not by assumption.

MR-1A launched on December~19, 1960, twenty-eight days later. It flew perfectly. Two hundred and thirty-five seconds of powered flight, 130 miles altitude, 235 miles downrange. The capsule separated, reentered, deployed its parachutes at the correct altitude over the correct ocean, and was recovered by helicopter. Every system worked. This time, success.

Wernher von Braun, watching MR-1A arc over the Atlantic, reportedly said something about the first one being a good lesson. Everyone nodded. It was. The best lessons always look stupid in retrospect because the knowledge they install is the kind you can never forget. Four inches. A quarter-inch difference in pin length. The entire future of American crewed spaceflight, briefly, hung on a race condition in a connector that nobody had tested at that specific rise height because nobody imagined the rocket would stop at four inches.

\begin{center}
$\bullet\quad\bullet\quad\bullet$
\end{center}

\textit{Il faut cultiver notre jardin}, Voltaire wrote at the end of \textit{Candide}. We must cultivate our garden. After everything --- the earthquakes, the shipwrecks, the Inquisition, the syphilis, the optimism --- the answer is to tend what's in front of you. Not because the world is the best of all possible worlds, but because the garden is the best of all possible responses to the world.

The garden at Launch Complex~5 was a rocket that needed a longer pin. We cultivated it. We made the pin a quarter inch longer. Then we launched again.

Four inches.

\vspace{2em}
\begin{center}
\textit{Dedicated to Voltaire (1694--1778)}\\[0.5em]
\textit{``Il faut cultiver notre jardin.''}\\
We must cultivate our garden.
\end{center}

\vspace{1em}
\noindent\rule{\textwidth}{0.4pt}
\begin{center}
\small NASA Mission Series v1.16 --- Mercury-Redstone 1\\
\small March 30, 2026
\end{center}

\end{document}
